

\section{Additional related works} \label{sec:additional_related_works}
\looseness-1\textbf{Interpretability:} Mechanistic interpretability works reverse-engineer the algorithms simulated by these models~\citep{elhage2021mathematical, olsson2022context, wang2022interpretability, nanda2023progress, chughtai2023toy, conmy2023towards}. These works study local patterns, e.g. activations and attention heads, to derive interpretable insights. 
Other works~\citep{weiss2021thinking,lindner2023tracr} use declarative programs to algorithmically describe transformer models. \citet{zhou2023algorithms} use these to explain task-specific length generalization of transformer models.


\textbf{Alternative Explanations for ICL:} Some works study ICL using a Bayesian framework. \citet{xie2022an} model pretraining data as a mixture of HMMs and cast ICL identifying one such component. \citet{hahn2023theory} later modeled language as a compositional grammar, and propose ICL as a composition of operations. \cite{zhang2023and,jiang2023latent,wang2023large,wies2023learnability} further strengthen this hypothesis by generalizing the underlying latent space. 
On the other hand, careful experiments in \citet{chan2022data} show that data distributional properties (e.g. Zipf's law) drive in-context learning in transformers.  

\begin{table*}[t]
  \centering
  \caption{Number of parameters of \simulator{}  for the forward, backward, and gradient update operations on various modules. For simplicity, we have ignored biases in the following computation. We set $S=4$, i.e. stack $4$ weights in each prefix embedding. We set $\simu{H}=12$ for \textsc{OPT-125M} and $\simu{H}=16$ for the other models, $\simu{D} = 4 \aux{D}$ for all the models, and $\simu{T} = \aux{T} + K$, with $\aux{T}=2048$ for \textsc{opt} models, and $K=\aux{D}/4$. $Q=4 Q_{ split} + 3 \simu{T} \simu{D} / \simu{H} $, where $Q_{ split } = \frac{1}{ \simu{H} }(  \simu{D} )^2 +  \simu{H} \simu{D}$,  denotes the number of parameters in a \simulator{} Linear Forward module (\cref{sec:exposition_linear_forward}). % $\aux{N}$ denotes the total number of parameters in the auxiliary model and needs to be included in the final \simulator{} size since they are stored in the Embedding matrix.
  }
  \label{tab:construction}
   %\mz{maybe add vanilla construction vs. construction after approximation? and make it clear how the size and flops increases with $h_d$.}
  \begin{tabular}{lcccc}
    \toprule
    & \multicolumn{4}{c}{Module Size} \\ \cmidrule(lr){2-5} 
    Module Name & Forward & Backward & Descent  & Total \\
    \midrule
    Linear layer  & $Q$ & $Q$ & $Q$  & $3Q$\\
    Layer norms & $Q$ & $Q + 2 \simu{D} \simu{H}$ & $Q$ & $3Q + 2 \simu{D} \simu{H}$\\
    Self-Attention  & $2Q$ & $2Q$ & $2Q$  & $6Q$\\
    Activation & $Q_{ split }$ & $2 \simu{D} \simu{H} $ & $0$ & $Q_{ split } + 2 \simu{D} \simu{H}$\\ \midrule
    Self-Attention block & $4Q$ &  $ 4Q+ 2\simu{D} \simu{H} $  & $ 4 Q $ & $ 12Q + 2\simu{D} \simu{H} $ \\
    Feed-forward block & $3Q+Q_{ split }$ & $ 3Q + 4\simu{D} \simu{H} $ & $ 3 Q $  & $ 9Q + 4\simu{D} \simu{H} $\\
    Transformer block & $7Q+Q_{ split }$ & $7Q + 6\simu{D} \simu{H}$ & $7Q$ & $ 21Q + 6\simu{D} \simu{H}+Q_{ split } $\\
    Transformer & $7QL + LQ_{ split }$ & $(7Q + 6 \simu{D} \simu{H})L$ & $7QL$ &  $ (21Q + 6\simu{D} \simu{H}+Q_{ split })L $ \\ % + $\aux{N}$\\ 
    \midrule
    \textsc{OPT-125m} & \textsc{0.4b} & \textsc{0.4b} & \textsc{0.4b} & \textsc{1.2b}       \\
    \textsc{OPT-350m} & \textsc{1.2b} & \textsc{1.1b} &  \textsc{1.1b} & \textsc{3.4b}       \\
    \textsc{OPT-1.3b} &  \textsc{3.7b}   &  \textsc{3.6b}   &    \textsc{3.5b} & \textsc{10.8b}  \\
    \textsc{OPT-2.7b} & \textsc{7.4b}   &  \textsc{7.2b}   &    \textsc{7.2b} & \textsc{21.8b}\\
    \bottomrule
  \end{tabular}
\end{table*}





\textbf{Transfer learning:} Our construction uses a pre-trained model to initialize a larger transformer, which is similar to several other more empirically oriented works~\citep{gong2019efficient,reddi2023efficient}. 


\section{Deferred defintions from main paper}\label{sec:deferred}

For simplicity of exposition, we showcase the definition on a single head self-attention layer (multi-head attention is in \cref{def:self-attn}).
%The self-attention in the auxiliary model is the same as in \simulator{} (\Cref{def:self-attn_construct_single}) without a position vector.
\begin{definition}[Auxiliary model softmax self-attention]
\label{def:self-attn_auxiliary_single}
    A self-attention layer with parameters $\{\mW_Q, \mW_K, \mW_V\}$ takes a sequence $\{ \vx_t \}_{t \le \aux{T}}$ and outputs a sequence $\{ \vy_t \}_{t \le \aux{T}}$, such that
    $$
    \vy_t = \sum_{j} a_{t, j} \vv_j, \qquad \text{with } a_{t, j} =\mathrm{softmax}(\mK \vq_t)_j, \quad 
    \vq_t = \mW_Q \vx_t, \quad \vk_t = \mW_K \vx_t, \quad \vv_t = \mW_V \vx_t,
    $$
    for all $t \le \aux{T}$, and $\mK \in \RR^{ \aux{T} \times  \aux{D} }$ defined with rows $\{ \vk_t \}_{t=1}^{\aux{T}}.$

\end{definition}




